\section{Repozytoria i uruchomienie}

\begin{table}[H]
\centering
\small
\begin{tabular}{lp{10.5cm}}
\toprule
repozytorium & zawartość \\
\midrule
\texttt{ryczek02/ticketbench} & repozytorium spinające: README, \texttt{ticketbench.sh}, runbook (\texttt{docs/runbook.html}), to sprawozdanie z~surowymi wynikami, submoduły pozostałych repozytoriów \\
\texttt{ticketbench-contracts} & pliki \texttt{.proto}, wygenerowany kod, model, generator danych, scenariusze, telemetria \\
\texttt{ticketbench-services} & \texttt{bench-runner}, \texttt{store-postgres}, \texttt{store-mariadb}, \texttt{store-mongo}, \texttt{store-valkey}, testy zgodności \\
\texttt{ticketbench-cli} & \texttt{tbench}: plany, zapis wyników, generator tabel i~wykresów \\
\texttt{ticketbench-infra} & Docker Compose, OTel/Jaeger/Prometheus/Grafana, macierze, Ansible, \texttt{bench.sh} \\
\bottomrule
\end{tabular}
\end{table}

Odtworzenie wszystkich wyników:
\begin{Verbatim}[frame=single,fontsize=\small]
git clone --recursive https://github.com/ryczek02/ticketbench
cd ticketbench
./ticketbench.sh all                # narzędzia, pobranie z GitHuba, stos, test dymny
./ticketbench.sh bench              # pełna macierz
./ticketbench.sh report             # tabele, wykresy, PDF
\end{Verbatim}

Wybrane polecenia CLI:
\begin{Verbatim}[frame=single,fontsize=\small]
tbench ping
tbench run --stores postgres,valkey --tickets 1M --scenarios R4,U6 --reps 10
tbench contend --clients 1,10,100,1000 --strategies all
tbench report results/<run> --out ../report/generated
\end{Verbatim}

Wszystkie liczby w~tabelach i~na wykresach pochodzą z~plików w~katalogu
\texttt{generated/}, które \texttt{tbench report} tworzy z~surowych wyników
w~\texttt{data/results}. CI repozytorium \texttt{ticketbench} generuje te pliki
ponownie i~sprawdza, czy są identyczne z~zatwierdzonymi. Szczegółowa instrukcja
uruchomienia jest w~runbooku \texttt{docs/runbook.html}.
